% v2.7.0 figure UID: FIG-P668-01
% Deterministic analytic Dirichlet densities on an N=24 barycentric mesh.
\begingroup
\tikzset{slfig-FIG-P668-01/.style={font=\fontsize{9.6pt}{11.5pt}\selectfont,
  every path/.append style={line cap=round,line join=round}}}
\pgfkeysifdefined{/tikz/alt}{}{\tikzset{alt/.code={}}}
\begin{figure}[htbp]
\centering
\begin{tikzpicture}[slfig-FIG-P668-01,slfig high,>=Stealth,
  every node/.style={font=\fontsize{9.6pt}{11.5pt}\selectfont,align=center},
  alt={对象—关系—结论：三个同尺度三分类单纯形均由给定Dirichlet密度公式逐网格计算；当全部形状分量为2时深色集中于中心且三面密度趋零，当全部为1时全域同色且密度恒为2，当全部为1/2时深色移向各面和顶点邻域且密度在边界发散。因此形状参数跨过1改变的是密度的边界行为，而不是单纯形本身。}]
\def\SLPN{24}
\def\SLPL{3.6}
\def\SLPH{3.117691454}
\def\SLPellmin{-1.224073476033}
\def\SLPellmax{2.777394473238}

% #1 selects Dir(2,2,2), Dir(1,1,1), or Dir(1/2,1/2,1/2).
% #2--#4 are a strictly interior barycentric sample point.
\newcommand{\SLPSetDensity}[4]{%
  \ifcase#1\relax
    \pgfmathsetmacro{\SLPdensity}{120*(#2)*(#3)*(#4)}%
  \or
    \pgfmathsetmacro{\SLPdensity}{2}%
  \or
    \pgfmathsetmacro{\SLPdensity}{1/(2*pi*sqrt((#2)*(#3)*(#4)))}%
  \fi
  \pgfmathsetmacro{\SLPell}{ln(\SLPdensity)}%
  \pgfmathsetmacro{\SLPfraction}{max(0,min(.999999,(\SLPell-\SLPellmin)/(\SLPellmax-\SLPellmin)))}%
  \pgfmathtruncatemacro{\SLPband}{floor(6*\SLPfraction)}%
  \ifcase\SLPband\relax
    \def\SLPshade{12}%
  \or\def\SLPshade{27}%
  \or\def\SLPshade{42}%
  \or\def\SLPshade{57}%
  \or\def\SLPshade{72}%
  \or\def\SLPshade{87}%
  \fi}

% Fixed shared log-density thresholds receive redundant line-style encoding.
\newcommand{\SLPSetContourStyle}[1]{%
  \ifcase#1\relax
    \def\SLPdash{densely dotted}%
  \or\def\SLPdash{densely dotted}%
  \or\def\SLPdash{dotted}%
  \or\def\SLPdash{densely dashed}%
  \or\def\SLPdash{dash dot}%
  \or\def\SLPdash{solid}%
  \fi}

\newcommand{\SLPFillPanel}[1]{%
  % Upward cells: centroid ((i+1/3)/N,(j+1/3)/N) is strictly interior.
  \foreach \SLPi in {0,...,23}{%
    \pgfmathtruncatemacro{\SLPjmax}{23-\SLPi}%
    \foreach \SLPj in {0,...,\SLPjmax}{%
      \SLPSetDensity{#1}
        {1-(\SLPi+.333333333333)/\SLPN-(\SLPj+.333333333333)/\SLPN}
        {(\SLPi+.333333333333)/\SLPN}{(\SLPj+.333333333333)/\SLPN}%
      \path[fill=SLBlue!\SLPshade!white,draw=SLBlue!\SLPshade!white,line width=.04pt]
        ({\SLPL*(\SLPi/\SLPN)+.5*\SLPL*(\SLPj/\SLPN)},{\SLPH*(\SLPj/\SLPN)})--
        ({\SLPL*((\SLPi+1)/\SLPN)+.5*\SLPL*(\SLPj/\SLPN)},{\SLPH*(\SLPj/\SLPN)})--
        ({\SLPL*(\SLPi/\SLPN)+.5*\SLPL*((\SLPj+1)/\SLPN)},{\SLPH*((\SLPj+1)/\SLPN)})--cycle;
    }%
  }%
  % Downward cells: centroid ((i+2/3)/N,(j+2/3)/N) is strictly interior.
  \foreach \SLPi in {0,...,22}{%
    \pgfmathtruncatemacro{\SLPjmax}{22-\SLPi}%
    \foreach \SLPj in {0,...,\SLPjmax}{%
      \SLPSetDensity{#1}
        {1-(\SLPi+.666666666667)/\SLPN-(\SLPj+.666666666667)/\SLPN}
        {(\SLPi+.666666666667)/\SLPN}{(\SLPj+.666666666667)/\SLPN}%
      \path[fill=SLBlue!\SLPshade!white,draw=SLBlue!\SLPshade!white,line width=.04pt]
        ({\SLPL*((\SLPi+1)/\SLPN)+.5*\SLPL*(\SLPj/\SLPN)},{\SLPH*(\SLPj/\SLPN)})--
        ({\SLPL*(\SLPi/\SLPN)+.5*\SLPL*((\SLPj+1)/\SLPN)},{\SLPH*((\SLPj+1)/\SLPN)})--
        ({\SLPL*((\SLPi+1)/\SLPN)+.5*\SLPL*((\SLPj+1)/\SLPN)},{\SLPH*((\SLPj+1)/\SLPN)})--cycle;
    }%
  }%
  % Draw only boundaries between formula-derived fixed bands, not hand-drawn contours.
  \foreach \SLPi in {0,...,23}{%
    \pgfmathtruncatemacro{\SLPjmax}{23-\SLPi}%
    \foreach \SLPj in {0,...,\SLPjmax}{%
      \pgfmathtruncatemacro{\SLPsum}{\SLPi+\SLPj}%
      \SLPSetDensity{#1}
        {1-(\SLPi+.333333333333)/\SLPN-(\SLPj+.333333333333)/\SLPN}
        {(\SLPi+.333333333333)/\SLPN}{(\SLPj+.333333333333)/\SLPN}%
      \edef\SLPbandA{\SLPband}%
      \ifnum\SLPsum<23\relax
        \SLPSetDensity{#1}
          {1-(\SLPi+.666666666667)/\SLPN-(\SLPj+.666666666667)/\SLPN}
          {(\SLPi+.666666666667)/\SLPN}{(\SLPj+.666666666667)/\SLPN}%
        \edef\SLPbandB{\SLPband}%
        \ifnum\SLPbandA=\SLPbandB\relax\else
          \pgfmathtruncatemacro{\SLPthreshold}{max(\SLPbandA,\SLPbandB)}%
          \SLPSetContourStyle{\SLPthreshold}%
          \draw[draw=SLRule,line width=.23pt,\SLPdash]
            ({\SLPL*((\SLPi+1)/\SLPN)+.5*\SLPL*(\SLPj/\SLPN)},{\SLPH*(\SLPj/\SLPN)})--
            ({\SLPL*(\SLPi/\SLPN)+.5*\SLPL*((\SLPj+1)/\SLPN)},{\SLPH*((\SLPj+1)/\SLPN)});
        \fi
      \fi
      \ifnum\SLPj>0\relax
        \SLPSetDensity{#1}
          {1-(\SLPi+.666666666667)/\SLPN-(\SLPj-.333333333333)/\SLPN}
          {(\SLPi+.666666666667)/\SLPN}{(\SLPj-.333333333333)/\SLPN}%
        \edef\SLPbandB{\SLPband}%
        \ifnum\SLPbandA=\SLPbandB\relax\else
          \pgfmathtruncatemacro{\SLPthreshold}{max(\SLPbandA,\SLPbandB)}%
          \SLPSetContourStyle{\SLPthreshold}%
          \draw[draw=SLRule,line width=.23pt,\SLPdash]
            ({\SLPL*(\SLPi/\SLPN)+.5*\SLPL*(\SLPj/\SLPN)},{\SLPH*(\SLPj/\SLPN)})--
            ({\SLPL*((\SLPi+1)/\SLPN)+.5*\SLPL*(\SLPj/\SLPN)},{\SLPH*(\SLPj/\SLPN)});
        \fi
      \fi
      \ifnum\SLPi>0\relax
        \SLPSetDensity{#1}
          {1-(\SLPi-.333333333333)/\SLPN-(\SLPj+.666666666667)/\SLPN}
          {(\SLPi-.333333333333)/\SLPN}{(\SLPj+.666666666667)/\SLPN}%
        \edef\SLPbandB{\SLPband}%
        \ifnum\SLPbandA=\SLPbandB\relax\else
          \pgfmathtruncatemacro{\SLPthreshold}{max(\SLPbandA,\SLPbandB)}%
          \SLPSetContourStyle{\SLPthreshold}%
          \draw[draw=SLRule,line width=.23pt,\SLPdash]
            ({\SLPL*(\SLPi/\SLPN)+.5*\SLPL*(\SLPj/\SLPN)},{\SLPH*(\SLPj/\SLPN)})--
            ({\SLPL*(\SLPi/\SLPN)+.5*\SLPL*((\SLPj+1)/\SLPN)},{\SLPH*((\SLPj+1)/\SLPN)});
        \fi
      \fi
    }%
  }}

\newcommand{\SLPPanel}[4]{%
  \begin{scope}[shift={(#2,0)}]
    \SLPFillPanel{#1}
    \draw[draw=SLRule,line width=.78pt] (0,0)--(\SLPL,0)--(.5*\SLPL,\SLPH)--cycle;
    \node[font=\fontsize{10.1pt}{12.1pt}\selectfont\bfseries,text=SLBlue]
      at (.5*\SLPL,3.58) {#3};
    \node[fill=white,fill opacity=.90,text opacity=1,rounded corners=1pt,
      inner xsep=2pt,inner ysep=1pt] at (.5*\SLPL,-.47) {#4};
    \node[anchor=south west,inner sep=1.5pt] at (0,0) {$(1,0,0)$};
    \node[anchor=south east,inner sep=1.5pt] at (\SLPL,0) {$(0,1,0)$};
    \node[anchor=north,inner sep=1.5pt] at (.5*\SLPL,\SLPH) {$(0,0,1)$};
    \ifcase#1\relax
      \draw[draw=SLRule,line width=.72pt] (.5*\SLPL,1.039230485) circle (2.8pt);
      \draw[draw=SLRule,line width=.72pt] ({.5*\SLPL-.07},1.039230485)--({.5*\SLPL+.07},1.039230485);
      \draw[draw=SLRule,line width=.72pt] (.5*\SLPL,.969230485)--(.5*\SLPL,1.109230485);
      \node[fill=white,rounded corners=1pt,inner sep=1.5pt] (SLPmode)
        at (.5*\SLPL,1.70) {唯一内点众数};
      \draw[-{Stealth[length=1.5mm]},draw=SLRule,line width=.52pt]
        (SLPmode.south) -- (.5*\SLPL,1.12);
      \node[fill=white,inner sep=1pt] at (.5*\SLPL,.14) {$p\to0$};
      \node[fill=white,inner sep=1pt,rotate=60] at (.53,1.06) {$p\to0$};
      \node[fill=white,inner sep=1pt,rotate=-60] at (3.07,1.06) {$p\to0$};
    \or
      \node[fill=white,rounded corners=1pt,inner sep=2pt,
        font=\fontsize{10.1pt}{12.1pt}\selectfont\bfseries] at (.5*\SLPL,1.30)
        {$p\equiv2$（全域均匀）};
    \or
      \node[fill=white,inner sep=1pt] at (.5*\SLPL,.14) {$p\to\infty$};
      \node[fill=white,inner sep=1pt,rotate=60] at (.53,1.06) {$p\to\infty$};
      \node[fill=white,inner sep=1pt,rotate=-60] at (3.07,1.06) {$p\to\infty$};
      \node[fill=white,rounded corners=1pt,inner sep=1.5pt] (SLPvertex) at (.5*\SLPL,2.12)
        {顶点邻域最强};
      \draw[-{Stealth[length=1.7mm]},draw=SLRule,line width=.55pt]
        (SLPvertex.north) -- ({.5*\SLPL-.20},2.76);
    \fi
  \end{scope}}

\SLPPanel{0}{0}{\textbf{(a)} 全部 $\alpha_i>1$\\$\alpha=(2,2,2)$}
  {$p=120\theta_1\theta_2\theta_3$}
\SLPPanel{1}{4.10}{\textbf{(b)} 全部 $\alpha_i=1$\\$\alpha=(1,1,1)$}
  {$p=2$}
\SLPPanel{2}{8.20}{\textbf{(c)} 全部 $\alpha_i<1$\\$\alpha=(1/2,1/2,1/2)$}
  {$p=(2\pi)^{-1}(\theta_1\theta_2\theta_3)^{-1/2}$}

\node[draw=SLRule,fill=white,rounded corners=2pt,text width=126mm,
  inner xsep=2.5mm,inner ysep=1.3mm] at (5.90,-1.08)
  {$\alpha_i>0,\quad \boldsymbol\theta\in\Delta_2^{\circ},\quad
    p(\boldsymbol\theta\mid\boldsymbol\alpha)=B(\boldsymbol\alpha)^{-1}
    \prod_{i=1}^3\theta_i^{\alpha_i-1}$；三组均值均为 $(1/3,1/3,1/3)$。};

\node[anchor=east] at (2.18,-1.79) {$\ell=\log p$};
\foreach \SLPk/\SLPshade in {0/12,1/27,2/42,3/57,4/72,5/87}{%
  \path[fill=SLBlue!\SLPshade!white]
    ({2.28+.68*\SLPk},-1.98) rectangle ({2.96+.68*\SLPk},-1.62);}
\draw[draw=SLRule,line width=.48pt] (2.28,-1.98) rectangle (6.36,-1.62);
\node[anchor=north] at (2.28,-2.00) {$-1.224$};
\node[anchor=north] at (6.36,-2.00) {$2.777$};
\node[anchor=west,align=left] at (6.62,-1.80)
  {共同色阶；五条色带边界为共同固定阈值\\
   $N=24$，每个小三角只在严格内点重心求值};
\node[anchor=north] at (5.90,-2.35)
  {$\varepsilon=0.01$ 只定义显示端点；真实边界仅按 $0/2/\infty$ 的解析极限标注。};
\end{tikzpicture}
\captionsetup{width=.96\linewidth}
\caption{三元Dirichlet的全部形状分量从大于1跨到等于1再到小于1时，密度由内部集中变为均匀，并在边界发散。}
\label{fig:V5-C05-dirichlet-shape-atlas}
\end{figure}
\endgroup
